\chapter{前言}
高等数学的计算建立在三个问题上：对象是什么、定理的条件是否满足、变形是否保留原问题。本册将极限、微分、积分和级数放进同一条逻辑链：用极限定义局部变化，用积分累积局部量，用微分方程描述变化规律，用级数构造和逼近函数。

全书按九章组织函数与极限、一元微积分、常微分方程、空间几何、多元微分、重积分、曲线曲面积分及无穷级数。每章由知识地图进入概念和定理，再用方法卡与例题展示计算、证明和几何解释，最后配备分层习题与方法速查表。

例题由基础到综合递进，多解训练着重比较代数、几何和分析中的不同入口。极限计算强调主项与余项，积分计算强调结构识别和换元推导，微分方程强调定义域与失解检查，级数强调收敛条件与端点处理。重要疑难点在警示盒及解答末尾说明。

全书包含596个例题、81道章末习题及解释图。习题附完整解答并提供双向链接；方法导航便于围绕同一数学结构组织复习。图注说明坐标、参数、定向和适用边界，阅读时应将图形直观与公式条件相互核对。

“考研风格”和“竞赛风格”表示题型与训练层次。不同考试科目的范围有所差异，可结合学习目标选择章节。高级专题区分完整证明、证明思路和标准结论，阅读时应先确认各项假设，再理解结论的适用范围。
